\newpage
\section{Derivadas de estabilidade}

% Secao 3 do roteiro: derivadas para o CG traseiro (aft.avl), sem trimagem,
% no ponto de projeto. Comandos ft, st e sb conforme as Tabelas 5 e 8 do
% enunciado. Lembretes de conversao para o MVO:
%   - derivadas de controle (delta_e, delta_a, delta_r, it) saem em 1/grau:
%     multiplicar por 180/pi;
%   - inverter o sinal de CYbeta, CYp, CYr, Cl_da, Cl_dr, Cn_da, Cn_dr;
%   - derivadas latero-direcionais de momento saem do sb, nao do st;
%   - inverter o sinal de zp.
% PREENCHER as tabelas apos as rodadas.

Os dados gerais da aeronave para a análise de estabilidade e controle,
extraídos do \texttt{designTool} no ponto de projeto, estão na
Tabela~\ref{tab:dados_gerais}. Os momentos de inércia foram calculados pela
função auxiliar da disciplina com a fração de combustível de 44\% que
reproduz o peso do ponto de projeto, com 100\% de carga paga, e são
referidos ao CG desse carregamento.

\begin{table}[H]
\centering
\caption{Informações gerais para análise de estabilidade e controle.}\label{tab:dados_gerais}
\renewcommand{\arraystretch}{1.25}
\begin{tabular}{lcl}
\toprule
Parâmetro & Valor & Fonte \\
\midrule
$S_{\mathrm{ref}}$ [m$^2$] & 368,833 & \texttt{designTool} \\
$c_{\mathrm{ref}}$ [m] & 6,8995 & \texttt{designTool} \\
$b_{\mathrm{ref}}$ [m] & 60,1518 & \texttt{designTool} \\
$m$ [kg] & 229.669,3 & \texttt{designTool} \\
$I_{xx}$ [kg$\cdot$m$^2$] & $1{,}508\times 10^{7}$ & \texttt{designTool} \\
$I_{yy}$ [kg$\cdot$m$^2$] & $4{,}284\times 10^{7}$ & \texttt{designTool} \\
$I_{zz}$ [kg$\cdot$m$^2$] & $5{,}578\times 10^{7}$ & \texttt{designTool} \\
$I_{xz}$ [kg$\cdot$m$^2$] & $1{,}209\times 10^{6}$ & \texttt{designTool} \\
$i_p$ [graus] & 0,0 & \texttt{designTool} \\
$x_p$ [m] & 18,0 & \texttt{designTool} \\
$z_p$ [m] & 3,0 & \texttt{designTool} (sinal invertido para o MVO) \\
$T_{\mathrm{max}}$ [N] & 1.026.000 & \texttt{designTool} (513.000 por motor, 2 motores) \\
$V$ [m/s] & 252,1 & ponto de projeto \\
$h$ [m] & 10.668 & ponto de projeto \\
\bottomrule
\end{tabular}
\end{table}

\subsection*{Resultados em $\alpha = 0$ e ajuste da polar}

% Tabela 6 do enunciado (CL0 e CM0 do comando ft em alpha = 0, it = 0,
% d2 = 0) e Tabela 7 (ajuste CD = CD0 + CDa*alpha + CDa2*alpha^2 sobre os
% dados da polar nao trimada do CG traseiro, item 5.b).

\begin{table}[H]
\centering
\caption{Resultados para $\alpha = 0$, $i_t = 0$ e ajuste quadrático da polar.}\label{tab:cl0_polar}
\renewcommand{\arraystretch}{1.25}
\begin{tabular}{lcl}
\toprule
Parâmetro & Valor & Fonte \\
\midrule
$C_{L0}$ & --- & AVL \texttt{ft} \\
$C_{M0}$ & --- & AVL \texttt{ft} \\
$C_{D0}$ & --- & ajuste do item 5.b \\
$C_{D\alpha}$ [1/rad] & --- & ajuste do item 5.b \\
$C_{D\alpha 2}$ [1/rad$^2$] & --- & ajuste do item 5.b \\
\bottomrule
\end{tabular}
\end{table}

\subsection*{Tabela completa de derivadas}

% Tabela 9 do enunciado, preenchida com ft, st e sb no ponto de projeto
% (a c 0.5053, it da trimagem do aft.avl, d2 = 0). Manter as conversoes de
% sinal e unidade indicadas nos comentarios do preambulo desta secao.
